Míra zaměstnanosti ve věkové skupině 20--64 let patří v~ČR dlouhodobě k~nejvyšším v~EU27, v~posledních letech přesahuje \SI{80}{\percent}.
Tato skutečnost je na první pohled příznivá, avšak v~kontextu teze představuje spíše varování: vysoká zaměstnanost při nízkém kolektivním pokrytí a podprůměrných mzdách (srov.~obr.~\ref{fig:net_income_ratio_timeline}) svědčí o~modelu „levné práce“, který nevytváří tlak na produktivitní inovace ani na mzdovou konvergenci.
Skandinávské země dosahují srovnatelné zaměstnanosti zároveň s~výrazně vyšší produktivitou a~mzdovou úrovní --- klíčový rozdíl leží v~institucionální architektuře trhu práce, zejména v~míře pokrytí kolektivními smlouvami a~výdajích na aktivní politiku zaměstnanosti.
Samotný indikátor zaměstnanosti proto nemůže sloužit jako argument proti reformě sociálního dialogu: nutno sledovat \emph{kvalitu} pracovních míst, nikoli pouze jejich počet.

Míra zaměstnanosti ve věkové skupině 20--64 let patří v~ČR dlouhodobě k~nejvyšším v~EU27, v~posledních letech přesahuje \SI{80}{\percent}.
Tato skutečnost je na první pohled příznivá, avšak v~kontextu teze představuje spíše varování: vysoká zaměstnanost při nízkém kolektivním pokrytí a podprůměrných mzdách (srov.~obr.~\ref{fig:net_income_ratio_timeline}) svědčí o~modelu „levné práce“, který nevytváří tlak na produktivitní inovace ani na mzdovou konvergenci.
Skandinávské země dosahují srovnatelné zaměstnanosti zároveň s~výrazně vyšší produktivitou a~mzdovou úrovní --- klíčový rozdíl leží v~institucionální architektuře trhu práce, zejména v~míře pokrytí kolektivními smlouvami a~výdajích na aktivní politiku zaměstnanosti.
Samotný indikátor zaměstnanosti proto nemůže sloužit jako argument proti reformě sociálního dialogu: nutno sledovat \emph{kvalitu} pracovních míst, nikoli pouze jejich počet.

Míra zaměstnanosti ve věkové skupině 20--64 let patří v~ČR dlouhodobě k~nejvyšším v~EU27, v~posledních letech přesahuje \SI{80}{\percent}.
Tato skutečnost je na první pohled příznivá, avšak v~kontextu teze představuje spíše varování: vysoká zaměstnanost při nízkém kolektivním pokrytí a podprůměrných mzdách (srov.~obr.~\ref{fig:net_income_ratio_timeline}) svědčí o~modelu „levné práce“, který nevytváří tlak na produktivitní inovace ani na mzdovou konvergenci.
Skandinávské země dosahují srovnatelné zaměstnanosti zároveň s~výrazně vyšší produktivitou a~mzdovou úrovní --- klíčový rozdíl leží v~institucionální architektuře trhu práce, zejména v~míře pokrytí kolektivními smlouvami a~výdajích na aktivní politiku zaměstnanosti.
Samotný indikátor zaměstnanosti proto nemůže sloužit jako argument proti reformě sociálního dialogu: nutno sledovat \emph{kvalitu} pracovních míst, nikoli pouze jejich počet.

Míra zaměstnanosti ve věkové skupině 20--64 let patří v~ČR dlouhodobě k~nejvyšším v~EU27, v~posledních letech přesahuje \SI{80}{\percent}.
Tato skutečnost je na první pohled příznivá, avšak v~kontextu teze představuje spíše varování: vysoká zaměstnanost při nízkém kolektivním pokrytí a podprůměrných mzdách (srov.~obr.~\ref{fig:net_income_ratio_timeline}) svědčí o~modelu „levné práce“, který nevytváří tlak na produktivitní inovace ani na mzdovou konvergenci.
Skandinávské země dosahují srovnatelné zaměstnanosti zároveň s~výrazně vyšší produktivitou a~mzdovou úrovní --- klíčový rozdíl leží v~institucionální architektuře trhu práce, zejména v~míře pokrytí kolektivními smlouvami a~výdajích na aktivní politiku zaměstnanosti.
Samotný indikátor zaměstnanosti proto nemůže sloužit jako argument proti reformě sociálního dialogu: nutno sledovat \emph{kvalitu} pracovních míst, nikoli pouze jejich počet.

\input{texparts/python/employment_rate_timeline}



